Sources : DGFiP, Statistiques n° 21, tableau 2 (IS brut et bénéfices, 2022) ; Insee Analyses n° 112 (septembre 2025), taux implicites 2022.\par Note : IS brut avant crédits d'impôt. Bénéfice fiscal : somme des bénéfices des sociétés bénéficiaires, après imputation des déficits reportés. Excédent net d'exploitation : ce que dégage l'activité, avant intérêts et impôt, au sens de la comptabilité nationale (Insee). Les deux sources ne mesurent pas exactement le même champ. Le taux réduit de 15 \% (bénéfice jusqu'à 42 500 \euro{}) explique l'essentiel du taux plus bas des microentreprises sur leur bénéfice.
